% from ega0-20-fr.tex
\oldpage[0\textsubscript{IV-1}]{150}
un idéal ouvert $\mathfrak{N}'\subset\mathfrak{N}$ de $B$ tels que
$\mathfrak{M}'B\subset\mathfrak{N}'$ et que l'image réciproque de $E$ par
l'homomorphisme canonique
$B/(\mathfrak{R}+\mathfrak{N}')\to B/(\mathfrak{R}+\mathfrak{N})$ soit
$(A/\mathfrak{M}')$-triviale.
\end{proof}

\begin{corollary}[20.7.9]
\label{0.20.7.9-fr}
Si la $A$-algèbre topologique $C=B/\mathfrak{R}$ est formellement lisse,
l'homomorphisme canonique $\delta_{C/B/A}$ est formellement inversible à
gauche.
\end{corollary}

\begin{env}[20.7.10]
\label{0.20.7.10-fr}
Dans \sref{0.20.6.1-fr}, lorsque $A$, $B$, $C$ sont des anneaux topologiques,
$u$ et $v$ des homomorphismes continus, on munit $\Upsilon_{C/B/A}$ de la
topologie induite par celle de $\Omega^1_{B/A}\otimes_B C$ ; les
homomorphismes \sref{0.20.6.4.2-fr} et \sref{0.20.6.4.4-fr} sont alors
continus pourvu qu'il en soit de même de ceux du diagramme
\sref{0.20.6.4.1-fr}. En outre, si, dans \sref{0.20.6.7-fr}, on suppose que
$L$ est un $C$-module discret annulé par un idéal ouvert de $C$, on déduit,
par passage à la limite inductive, un $C$-homomorphisme canonique
\begin{equation}
\tag{20.7.10.1}
\label{0.20.7.10.1-fr}
  \operatorname{Exalcotop}_{B/A}(C,L)\to
  \operatorname{Hom.\ cont}_C(\Upsilon_{C/B/A},L)
\end{equation}
compte tenu de \sref{0.18.5.3.1-fr} : pour tout idéal ouvert
$\mathfrak{M}$ de $A$, tout idéal ouvert $\mathfrak{N}$ de $B$ tel que
$\mathfrak{M}B\subset\mathfrak{N}$, tout idéal ouvert $\mathfrak{P}$ de $C$
tel que $\mathfrak{N}C\subset\mathfrak{P}$ et que $\mathfrak{P}.L=0$, toute
$(B/\mathfrak{N})$-extension $E$ de $C/\mathfrak{P}$ par $L$ qui est
$(A/\mathfrak{M})$-triviale provient de la donnée d'un $C$-homomorphisme
continu $f$ de $\Omega^1_{B/A}\otimes_B C$ dans $L$, et l'homomorphisme
\sref{0.20.7.10.1-fr} fait correspondre à l'image dans
$\operatorname{Exalcotop}_{B/A}(C,L)$ de la classe de $E$, la restriction de
$f$ à $\Upsilon_{C/B/A}$, dit \emph{homomorphisme caractéristique de $E$} et
noté encore $\chi_E$.
\end{env}

\begin{proposition}[20.7.11]
\label{0.20.7.11-fr}
Supposons que la topologie de $C$ soit telle que le carré d'un idéal ouvert
soit ouvert. Si $C$ est une $A$-algèbre topologique formellement lisse et si
$\Omega^1_{C/B}$ est un $C$-module formellement projectif, on a un
isomorphisme canonique
\begin{equation}
\tag{20.7.11.1}
\label{0.20.7.11.1-fr}
  \operatorname{Exalcotop}_B(C,L)\xrightarrow{\sim}
  \operatorname{Hom.\ cont}_C(\Upsilon_{C/B/A},L)
\end{equation}
pour tout $C$-module discret $L$ annulé par un idéal ouvert de $C$.
\end{proposition}

\begin{proof}
On a en effet, dans la suite exacte \sref{0.20.3.7.1-fr},
$\operatorname{Exalcotop}_A(C,L)=0$ \sref{0.19.4.4-fr}, donc
$\operatorname{Exalcotop}_B(C,L)=\operatorname{Exalcotop}_{B/A}(C,L)$ et
l'homomorphisme \sref{0.20.7.8.1-fr} n'est autre que
\sref{0.20.7.10.1-fr} ; le fait qu'il est bijectif se déduit de
\sref{0.20.6.13-fr} par passage à la limite inductive, compte tenu de ce que
la topologie de $\Omega^1_{C/B}$ est alors déduite de celle de $C$
\sref{0.20.4.5-fr} et de \sref{0.19.2.4-fr}.
\end{proof}

\begin{env}[20.7.12]
\label{0.20.7.12-fr}
Dans \sref{0.20.6.14-fr} on munit encore $\Upsilon^C_{B/A/\Lambda}$ de la
topologie induite par celle de $\Omega^1_{A/\Lambda}\otimes_A C$ et alors
l'homomorphisme \sref{0.20.6.14.6-fr} est continu, lorsque les anneaux
considérés sont topologiques et les homomorphismes d'anneaux continus.
\end{env}

\begin{env}[20.7.13]
\label{0.20.7.13-fr}
Si, dans \sref{0.20.6.23-fr}, on suppose que les anneaux $\Lambda$, $A$, $B$,
$C$ sont topologiques, les homomorphismes $s$, $u$, $v$ continus et $L$ un
$C$-module discret annulé par un idéal ouvert de $C$, alors on peut passer à
la limite inductive comme dans \sref{0.20.3.6-fr}, et on a une suite exacte
\begin{equation}
\tag{20.7.13.1}
\label{0.20.7.13.1-fr}
\begin{split}
  0\to\operatorname{D\acute{e}r.\ cont}_B(C,L)\to
  \operatorname{D\acute{e}r.\ cont}_A(C,L)\to
  \operatorname{D\acute{e}r.\ cont}_A(B,L)\to \\
  \to\operatorname{Exalcotop}_{B/\Lambda}(C,L)\to
  \operatorname{Exalcotop}_{A/\Lambda}(C,L)\to
  \operatorname{Exalcotop}_{A/\Lambda}(B,L)\to0.
\end{split}
\end{equation}
\end{env}

\begin{env}[20.7.14]
\label{0.20.7.14-fr}
Soient $A$ un anneau topologique, $B$ une $A$-algèbre topologique,
$\mathfrak{M}'$ un idéal ouvert de $A$, $\mathfrak{N}'$ un idéal ouvert de
$B$ tels que $\mathfrak{M}'B\subset\mathfrak{N}'$ ; posons
$A'=A/\mathfrak{M}'$,
